\begin{abstract}
Recent advancements in multi-task model merging offer a promising, training-free alternative for parameter integration, yet current state-of-the-art adaptive approaches rely heavily on heuristic test-time adaptation via unsupervised entropy minimization. While empirically functional, this paradigm lacks joint multi-task loss minimization guarantees, suffers from high computational overhead, and is highly prone to transductive overfitting and sacrificial task bias. To address these theoretical limitations, we introduce Curvature-Aware Analytical Model Merging (ACM), a mathematically rigorous framework that reformulates parameter fusion as a quadratic optimization problem based on local second-order Taylor expansions around task expert minima. By projecting the high-dimensional parameter space onto the low-dimensional subspace of the expert task vectors, ACM computes the full, non-diagonal, cross-parameter Hessian curvature with zero diagonal approximation. This yields a provably optimal, closed-form analytical solution for layer-wise merging coefficients that is entirely training-free and solved in a single step. To guarantee complete scientific integrity, we select hyperparameters strictly via an unsupervised calibration validation split heuristic, eliminating test-set leakage. Extensive simulation sweeps across 30 seeds show that ACM outperforms existing methods by up to +8.11\% in coupled non-convex landscapes. Furthermore, physical validation on Vision Transformers confirms that our proposed leakage-free Global-Normalized ACM (ACM-GlobalNorm) achieves a Joint Average accuracy of 57.76\%, significantly outperforming both diagonal Fisher Merging (56.03\%) and polynomial test-time adaptation (38.96\%) without any test-time optimization. While standard Task Arithmetic achieves 60.72\% when exhaustively tuned (scale 0.4) on the test set, we analyze the local quadratic approximation's behavior on highly non-convex landscapes, providing deep theoretical insights into curvature-aware weight consolidation.
\end{abstract}